Acrónimos usados en el presente trabajo de investigación:
\begin{acronym}[ABCDEFGHIJK]
%A
\acro{AF}{alta frecuencia}
%B
\acro{BF}{baja frecuencia}

%c
\acro{CoS}{concepto de solución}
%\ac{API}
%D
\acro{DC}{\textit{Direct Current}} %\ac{DC}
%E
\acro{ETAP}{\textit{Electrical Transient and Analysis Program}}
%F
\acro{FEM}{Fuerza contraelectromotriz}
\acro{FPGA}{\textit{Field-Programmable Gate Array}}
%G
\acro{GUI}{\textit{Graphical User Interface}}
%H
\acro{HW}{Hardware} %\ac{HW}
%M
\acro{MEE}{Modelo de Espacio de Estados}
\acro{MIMO}{\textit{Multiple Input Multiple Output}}
\acro{MCB}{\textit{Motor Control Blockset}}
\acro{MQTT}{\textit{Message Queuing Telemetry Transport}}
%N
\acro{TRL}{\textit{Technological Readiness Level}}
%P
\acro{PC}{\textit{Personal Computer}}
\acro{PCB}{\textit{Printed Circuit Board}}
\acro{PUCP}{Pontificia Universidad Católica del Perú}

%S
\acro{SW}{Software}
\acro{SWA}{Software A}
\acro{SWB}{Software B}
\acro{SISO}{\textit{Single Input Single Output}}
\acro{SBC}{\textit{Single Board Computer}}
%T
\acro{TFC1}{Trabajo de Fin de Carrera 1}
\acro{TFC2}{Trabajo de Fin de Carrera 2}
\end{acronym}